\section{Proof of Theorem 1.2}

From Lemma \ref{lem3.5}, we see that any critical point of $I_\mu|_{S_r(c)}$ belongs 
to $\mathcal{P}_c$. Consequently, the properties of the manifold $\mathcal{P}_c$ have  
relation to the mini-max structure of $I_\mu|_{ S_{r}(c)}$. For $u\in  S_{r}(c)$ and 
$\theta\in \mathbb{R}$, we introduce the transformation:
\[ %\label{5.1}
(\theta \star u)(x):=e^{\frac{3\theta}{2}}u(e^\theta x), \quad
x\in \mathbb{R}^3,\;\theta \in \mathbb{R}.
\]
It is easy to check that the dilations preserve the $L^2$-norm such that
$\theta \star u\in  S_{r}(c)$, by direct calculation, one has
\begin{align*}%\label{5.2}
I(u,\theta)
&= I_\mu((\theta \star u))\\
&:= \frac{a}{2}e^{2s\theta}\int_{\mathbb{R}^3}|(-\Delta)^{s/2}u|^2dx
+\frac{b}{4}e^{4s\theta}\Big(\int_{\mathbb{R}^3}|(-\Delta)^{s/2}u|^2dx\Big)^2\\
&\quad +\frac{1}{4}e^{(3-2s)\theta}\int_{\mathbb{R}^3}\phi^s_u u^2dx-\frac{\mu}{p}e^{\frac{3(p-2)}{2}\theta}\int_{\mathbb{R}^3}|u|^pdx-
\frac{1}{2^*_s}e^{\frac{3(2^*_s-2)}{2}\theta}\int_{\mathbb{R}^3}|u|^{2^*_s}dx.
\end{align*}

\begin{lemma}\label{lem5.1}
Let $u\in  S_{r}(c)$, then
\begin{itemize}
\item[(i)] $\int_{\mathbb{R}^3}|(-\Delta)^{s/2}(\theta \star u)|^2dx\to 0$ and 
$I_\mu((\theta \star u))\to 0$ as $\theta \to -\infty$;

\item[(ii)] $\int_{\mathbb{R}^3}|(-\Delta)^{s/2}(\theta \star u)|^2dx\to +\infty$ and 
$I_\mu((\theta \star u))\to -\infty$ as $\theta \to +\infty$.
\end{itemize}
\end{lemma}

\begin{proof}
 A direct computation shows that
 \begin{gather*} %\label{5.3}
\int_{\mathbb{R}^3}|(-\Delta)^{s/2}(\theta \star u)|^2dx=e^{2s\theta}
\int_{\mathbb{R}^3}|(-\Delta)^{s/2}u|^2dx, \\
\label{5.4}
\int_{\mathbb{R}^3}|(-\Delta)^{s/2}(\theta \star u)|^2dx\to 0\quad
\text{as }\theta \to -\infty, \\
\label{5.41}
\int_{\mathbb{R}^3}|(-\Delta)^{s/2}(\theta \star u)|^2dx\to +\infty\quad
\text{as }\theta \to +\infty.
\end{gather*}
Notice that
\begin{align*}%\label{5.2}
I_\mu((\theta \star u))
&:= \frac{a}{2}e^{2s\theta}\int_{\mathbb{R}^3}|(-\Delta)^{s/2}u|^2dx
+\frac{b}{4}e^{4s\theta}\Big(\int_{\mathbb{R}^3}|(-\Delta)^{s/2}u|^2dx\Big)^2\\
&\quad +\frac{1}{4}e^{(3-2s)\theta}\int_{\mathbb{R}^3}\phi^s_u u^2dx
 -\frac{\mu}{p}e^{\frac{3(p-2)}{2}\theta}\int_{\mathbb{R}^3}|u|^pdx
 -\frac{1}{2^*_s}e^{\frac{3(2^*_s-2)}{2}\theta}\int_{\mathbb{R}^3}|u|^{2^*_s}dx,
\end{align*}
by $\frac{3(2^*_s-2)}{2}>\frac{3(p-2)}{2}>4s>2s>3-2s$, we infer that
$$
 I_\mu((\theta \star u))\to 0 \quad\text{as }\theta \to -\infty,
 \quad I_\mu((\theta \star u))\to -\infty ~as ~\theta \to +\infty.
$$
This completes the proof.
\end{proof}

\begin{lemma}\label{lem5.2} 
There exist $K=K_c>0$ and $\tilde{c}>0$ such that for all $0<c<\tilde{c}$,
\begin{equation}\label{5.3}
0<\sup_{u\in\mathcal{A}_c}I_\mu(u)<\inf_{u\in\mathcal{B}_c}I_\mu(u),
\end{equation}
